% The core slide of the block: why value-based offline RL breaks.
% Deliberate contrast with the cliff: this failure needs no dynamics model.

% Goal/lead-in: WHY does running Q-learning on the fixed dataset break?
\begin{frame}{The Core Failure: Extrapolation Error}
\begin{itemize}
    \item Stitching needs \textbf{Q-learning}, so just run it on the fixed
          dataset. Look at its update target, which bootstraps from a $\max$ over
          actions:
          \[ y = r + \gamma\,\max_{a'} Q_\theta(s', a') \]
    \item That $\max$ ranges over \textbf{all} actions, including ones the
          dataset \emph{never tried} (\emph{out-of-distribution}, or \textbf{OOD}).
          Why is that fatal?
\end{itemize}
\end{frame}

\begin{frame}{The Core Failure: Extrapolation Error}
\definecolor{oteal}{RGB}{0,141,160}
\definecolor{ogreen}{RGB}{106,168,79}
\definecolor{ored}{RGB}{158,28,38}
\definecolor{oblue}{RGB}{16,53,103}
\centering
\begin{tikzpicture}[font=\scriptsize, >={Stealth[length=2.4mm]}]
    \draw[->] (0,0) -- (8.0,0) node[below left]{$a$};
    \draw[->] (0,0) -- (0,3.3);
    \node[rotate=90] at (-0.35,1.6) {$Q(s,\cdot)$};
    \fill[oteal!12] (0.8,0.02) rectangle (3.6,3.0);
    \node[oteal!60!black] at (2.2,2.8) {data support};
    \draw[ogreen!70!black, thick] plot[smooth] coordinates
      {(0.3,0.9) (1.2,1.5) (2.2,1.9) (3.2,1.7) (4.5,1.2) (6.0,0.85) (7.4,0.6)};
    \node[ogreen!60!black, font=\scriptsize] at (7.3,0.95) {true $Q$};
    \draw[ored, very thick, dashed] plot[smooth] coordinates
      {(0.3,1.0) (1.2,1.45) (2.2,1.95) (3.2,1.75) (4.4,1.9) (5.6,2.5) (7.0,3.05)};
    \node[ored, font=\scriptsize] at (7.45,2.75) {learned $\hat Q$};
    \foreach \x/\y in {1.0/1.42, 1.5/1.66, 1.9/1.84, 2.4/1.92, 2.8/1.85, 3.3/1.66}
      \fill[oblue] (\x,\y) circle (1.6pt);
    \node[ored, font=\scriptsize\bfseries] (ph) at (5.2,3.15) {phantom peak};
    \draw[->, ored] (ph.east) -- (6.85,3.0);
    \draw[->, ored, thick] (7.0,-0.5) -- (7.0,-0.08);
    \node[ored, font=\scriptsize] at (5.7,-0.5)
      {$\arg\max_a \hat Q$: the policy goes here};
\end{tikzpicture}
\begin{itemize}
    \item $Q_\theta$ is a neural net: feed it \emph{any} action and it returns a
          number, even where there is \textbf{no data}; there the value is not
          learned, only \emph{extrapolated} (``invented'')
    \item The $\max$ seeks the \emph{largest} invented value $\Rightarrow$
          \textbf{overestimation}; the policy $\arg\max_a \hat Q$ chases that
          \emph{phantom}, and bootstrapping spreads the inflated value onward
\end{itemize}
\end{frame}

\begin{frame}{Why Online RL Never Noticed}
\begin{itemize}
    \item Online, an overestimated action gets \emph{tried}, disappoints, and is
          corrected; overestimation is self-erasing
\end{itemize}
\end{frame}

\begin{frame}{Why Online RL Never Noticed}
\begin{itemize}
    \item Online, an overestimated action gets \emph{tried}, disappoints, and is
          corrected; overestimation is self-erasing
    \item Offline, the policy can chase the phantom forever: no data will ever
          contradict it
    \item More gradient steps make it \textbf{worse}, not better: we optimize
          against the fantasy (Fujimoto et al., 2019: \emph{extrapolation error})
\end{itemize}
\vspace{0.8em}
\begin{center}
\textcolor{red}{Online, hype meets reality. Offline, it never does.}
\end{center}
\end{frame}
